% !TEX root = ../../ctfp-print.tex

\lettrine[lhang=0.17]{Y}{ou can get} 通过研究各种实例可以获得对范畴的真正理解。范畴呈现出各种形态和规模，常常出现在意想不到的地方。我们将从最简单的例子开始。

\section{无对象}

最平凡的范畴包含零个对象（object）以及相应的零个态射（morphism）。单独来看这是一个非常乏味的范畴，但在其他范畴的上下文中可能具有重要意义，例如在全体范畴构成的范畴中（是的，确实存在这样的范畴）。如果你认为空集的概念是有意义的，那么为什么不能有空范畴呢？

\section{简单图}

你可以通过用箭头连接对象来构建范畴。可以想象从任意有向图出发，通过添加更多箭头将其转化为范畴。具体来说：首先，在每个节点处添加恒等箭头；接着，对于任意两个满足一个箭头的终点与另一个箭头的起点重合的箭头（换句话说，任意两个\newterm{可组合（composable）}的箭头），添加一个新的箭头作为它们的合成箭头。每当添加新箭头时，都必须考虑它与任何其他箭头（除恒等箭头外）以及自身的合成运算。这个过程通常会产生无限多个箭头，但这并不影响其有效性。

另一种看待此过程的方式是：你创建了一个范畴，该范畴的每个图节点对应一个对象，而所有可能的\newterm{可组合图边的链（chain）}构成态射集合（甚至可以将长度为零的链视为恒等态射的特殊情况）。

这种范畴被称为由给定图生成的\newterm{自由范畴（free category）}。这是"自由构造"的一个典型示例——通过向给定结构扩展最少必要元素以满足其公理体系（此处即范畴的三条定律）的过程。我们将在后续章节看到更多此类构造实例。

\section{序}

现在换一个完全不同的角度！考虑这样一个范畴：其中态射是对象之间的关系——即「小于等于」的关系。让我们验证它是否确实构成一个范畴。我们是否有恒等态射？每个对象都满足 $a \leqslant a$：检查通过！是否有合成法则？若 $a \leqslant b$ 且 $b \leqslant c$，则必有 $a \leqslant c$：检查通过！合成是否满足结合律？检查通过！这种具有如此关系的集合称为\newterm{预序（preorder）}，因此预序确实是一个范畴。

我们可以引入更强的关系，要求满足额外条件：若 $a \leqslant b$ 且 $b \leqslant a$，则 $a$ 必须等于 $b$。这被称为\newterm{偏序（partial order）}。

最后，你还可以施加更强的限制条件：任意两个对象之间必然存在单向关系，这样的结构称为\newterm{线序（linear order）}或\newterm{全序（total order）}。

让我们用范畴论特征刻画这些有序集合。预序是一个对任意对象 $a,b$ 至多存在一个态射的范畴。这类范畴的别称是「薄范畴（thin category）」，因此预序即薄范畴。

在范畴 $\cat{C}$ 中，从对象 $a$ 到对象 $b$ 的态射集称为\newterm{同态集（hom-set）}，记作 $\cat{C}(a, b)$（有时写作 $\mathbf{Hom}_{\cat{C}}(a, b)$）。因此预序中的每个同态集要么为空集，要么为单元素集。这包含 $\cat{C}(a, a)$ 这个特例——任何预序中从 $a$ 到 $a$ 的态射集必为仅含恒等态射的单元素集。然而，预序允许存在环路，而偏序则禁止环路。

能够识别预序、偏序和全序结构非常重要，因为这与排序算法密切相关。快速排序、冒泡排序、归并排序等算法仅能在全序结构上正确运行。偏序结构则可通过拓扑排序进行处理。

\section{作为集合的幺半群}

幺半群（monoid）是一个异常简单却极其强大的概念。它背后隐藏着基础算术的本质：加法和乘法都构成了幺半群。幺半群在编程中无处不在——它们以字符串、列表、可折叠数据结构的形式出现，在并发编程中表现为未来值（future），在响应式函数编程中体现为事件等等。

传统上，幺半群被定义为带有二元运算的集合。该运算只需满足两个条件：其一是结合律；其二是存在一个特殊元素，相对于该运算具有单位元（neutral element）的性质。

例如，包含零的自然数在加法下构成幺半群。结合律意味着：
$$(a + b) + c = a + (b + c)$$
（换句话说，进行加法运算时我们可以省略括号）

加法的单位元是零，因为：
$$0 + a = a$$
且
$$a + 0 = a$$
第二个方程实际上是冗余的，因为加法具有交换性 $(a + b = b + a)$，但交换性并不属于幺半群定义的一部分。例如，字符串拼接不具有交换性，但它仍然构成幺半群。顺便说一句，字符串拼接的单位元是空字符串，将其连接到任意字符串两侧都不会改变原字符串。

在Haskell中我们可以定义一个表示幺半群的类型类——对于某个类型，若存在名为 \code{mempty} 的单位元和名为 \code{mappend} 的二元运算，则满足该类型类：

\src{snippet01}
对于双参数函数的类型签名 \code{m -> m -> m}，初看可能有些奇怪，
但在讨论柯里化（currying）之后就会变得清晰。你可以用两种基本方式解读多重箭头的类型签名：将其视为多个参数的函数，最右侧类型为返回类型；或者视为接受单个参数（最左侧参数）并返回函数的函数。通过添加括号可以强调后一种解读方式（虽然这些括号是冗余的，因为箭头是右结合的），即写作：\code{m -> (m -> m)}。
我们稍后会回到这种解读方式。

需要注意的是，在Haskell中无法直接表达 \code{mempty} 和 \code{mappend} 的幺半群性质（即 \code{mempty} 的单位元性质和 \code{mappend} 的结合律）。程序员有责任确保这些性质得到满足。

Haskell的类型类不像C++类那样侵入性强。当定义新类型时，无需预先指定其类型类。你可以推迟决策，直到很晚才声明某个类型是某个类型类的实例。举例来说，让我们通过提供 \code{mempty} 和 \code{mappend} 的实现来声明 \code{String} 是幺半群（实际上标准预置库已经为你完成了这一点）：

\src{snippet02}
这里我们复用了列表拼接运算符 \code{(++)}，因为 \code{String} 本质上就是字符列表。

关于Haskell语法的一个说明：任何中缀运算符都可以通过添加括号转换为双参数函数。给定两个字符串，你可以在它们之间插入 \code{++} 进行拼接：

\begin{snip}{haskell}
"Hello " ++ "world!"
\end{snip}
或者将它们作为两个参数传递给带括号的 \code{(++)} 函数：

\begin{snip}{haskell}
(++) "Hello " "world!"
\end{snip}
请注意函数参数既不使用逗号分隔，也不包裹在括号中。（这可能是学习Haskell时最难适应的部分）

值得强调的是，Haskell允许你直接表达函数的相等性，例如：

\begin{snip}{haskell}
mappend = (++)
\end{snip}
从概念上讲，这不同于表达函数产生的值的相等性，后者形式如下：

\begin{snip}{haskell}
mappend s1 s2 = (++) s1 s2
\end{snip}
前者对应范畴 $\Hask$ 中态射（或忽略底部值时的 $\Set$ 范畴）的等价关系。这类方程不仅更为简洁，而且常常可以推广到其他范畴。后者被称为\newterm{外延相等性}，表明对于任意两个输入字符串，\code{mappend} 和 \code{(++)} 的输出结果相同。由于参数值有时被称为\newterm{点}（正如在某点 $x$ 处的函数值 $f(x)$），这种相等性也称为逐点（point-wise）相等性。不显式指定参数的函数相等性则称为\newterm{无点式}（point-free）。（顺便提一下，无点式方程经常涉及函数组合，其符号表示是一个点，这对初学者来说可能会有点困惑。）

在C++中声明幺半群最接近的方式是使用C++20标准的概念（concept）特性：

\begin{snip}{cpp}
template<class T>
struct mempty;

template<class T>
T mappend(T, T) = delete;

template<class M>
concept Monoid = requires (M m) {
    { mempty<M>::value() } -> std::same_as<M>;
    { mappend(m, m) } -> std::same_as<M>;
};
\end{snip}
第一个定义是用于保存每个特化类型的单位元的结构体。

关键字 \code{delete} 表示没有默认定义：必须针对每个具体类型单独指定。
同理，\code{mappend} 也没有默认实现。

概念 \code{Monoid} 用于测试给定类型 \code{M} 是否存在合适的 \code{mempty} 和 \code{mappend} 定义。

可以通过提供适当的特化和重载来实例化Monoid概念：

\begin{snip}{cpp}
template<>
struct mempty<std::string> {
    static std::string value() { return ""; }
};

template<>
std::string mappend(std::string s1, std::string s2) {
    return s1 + s2;
}
\end{snip}

\section{作为范畴的幺半群}

这是我们从集合元素的角度出发对幺半群的「熟悉」定义。但在范畴论中，我们总是试图摆脱集合及其元素的视角，转而采用对象与态射的观点。让我们稍作视角转换，把二元运算看作是在集合内部「移动」或「位移」的操作。

例如，存在一个向每个自然数添加5的操作：它将0映射到5，1映射到6，2映射到7，依此类推。这是定义在自然数集合上的一个函数——这很好，我们得到了函数与集合的概念。一般地，对于任意数$n$都存在对应的加法函数（即$n$的「加法器」）。

这些加法器如何组合？添加5的函数与添加7的函数组合后，就得到一个添加12的函数。因此加法器的组合规则可以对应到加法运算规则。这也很妙：我们可以用函数组合替代加法运算。

更进一步，还存在中性元（零）对应的加法器。添加零不会引起任何位移，这就是自然数集合上的恒等函数。

与其给出传统的加法规则，不如给出加法器的组合规则，这样并不会丢失任何信息。注意加法器的组合具有结合律（因为函数组合本身具有结合律），并且存在对应恒等函数的零加法器。

敏锐的读者可能已经注意到，整数到加法器的映射源自对类型签名\code{mappend}的第二种解读：\code{m -> (m -> m)}。这个签名表明\code{mappend}将幺半群集合中的元素映射到作用于该集合的函数。

现在请暂时忘记你面对的是自然数集合，而是将其视为单一对象——一个包含若干态射（加法器）的「斑点」。幺半群本质上是一个单对象范畴。事实上，「monoid」一词源自希腊语\emph{mono}（意为「单一」）。每个幺半群都可以描述为具有适当组合规则的单对象范畴。

\begin{figure}[H]
  \centering
  \includegraphics[width=0.35\textwidth]{images/monoid.jpg}
\end{figure}

\noindent
字符串连接是一个有趣案例，因为我们既可以选择定义右附加器（right appender）也可以选择左附加器（left appender，或称\emph{prependers}）。这两种模型的组合表互为镜像反转。你可以自行验证：在「foo」后追加「bar」相当于在「bar」前预置「foo」后再进行预置操作。

或许你会疑问：每个范畴论意义上的幺半群（即单对象范畴）是否都能唯一确定一个带二元运算的集合？事实证明我们总能从单对象范畴中提取出一个集合——就是我们的例子中的那些加法器构成的集合。换句话说，这个集合就是范畴$\cat{M}$中单对象$m$的同态集合$\cat{M}(m, m)$。我们可以在该集合上定义二元运算：两个集合元素的单子积对应它们代表的态射的组合。若给出$\cat{M}(m, m)$中对应$f$和$g$的两个元素，它们的乘积对应组合态射$f \circ g$。这种组合必然存在，因为这些态射的源与目标是同一对象；其结合律由范畴的基本规则保证；而恒等态射就是该乘法的单位元。因此我们总能从范畴论幺半群恢复出集合论幺半群。可以说两者在本质上是同一概念。

\begin{figure}[H]
  \centering
  \includegraphics[width=0.4\textwidth]{images/monoidhomset.jpg}
  \caption{作为态射和作为集合点的幺半群同态集合}
\end{figure}

\noindent
留给数学家的一个细节是：态射未必构成集合。在范畴世界中存在比集合更大的结构。若任意两个对象间的态射构成集合，则该范畴称为局部小范畴（locally small）。如前所承诺，本文将忽略这类细微差别，但为严谨起见仍需提及。

范畴论中许多有趣的现像源于这样一个事实：同态集合中的元素既可以视为遵循组合规则的态射，又可以视为集合中的点。在此例中，范畴$\cat{M}$中的态射组合就转化为集合$\cat{M}(m, m)$中的单子积。

\section{挑战}

\begin{enumerate}
  \tightlist
  \item
        从以下对象生成一个自由范畴：

        \begin{enumerate}
          \tightlist
          \item
                包含一个节点和零条边的图
          \item
                包含一个节点和一条（有向）边的图（提示：该边可以与自身复合）
          \item
                包含两个节点和一条单向箭头的图
          \item
                包含一个节点和26条用英文字母标记的箭头（a, b, c \ldots{} z）的图
        \end{enumerate}
  \item
        以下分别是哪种序结构？

        \begin{enumerate}
          \tightlist
          \item
                在包含关系下的集合族：当集合$A$的所有元素都属于集合$B$时，称$A$包含于$B$
          \item
                C++类型系统中如下子类型关系：当指向\code{T1}的指针可作为期望指向\code{T2}的参数传递给函数且不触发编译错误时，称\code{T1}是\code{T2}的子类型
        \end{enumerate}
  \item
        已知\code{Bool}是由\code{True}和\code{False}构成的二元集合，证明其相对于运算符\code{\&\&}（逻辑与）和\code{||}（逻辑或）分别构成两个（集合论意义上的）幺半群。
  \item
        将带有AND运算符的\code{Bool}幺半群表示为范畴：列出所有态射及其合成规则。
  \item
        将模3加法表示为幺半群范畴。
\end{enumerate}